% ============================================================================
% 中国古代史 · 下卷（7-9年级）写作模板
% ============================================================================

\documentclass[a4paper,12pt]{ctexart}
\usepackage{xcolor}
\usepackage{graphicx}
\usepackage[top=2cm, bottom=2cm, left=2cm, right=2cm]{geometry}
\usepackage{array}
\usepackage{multirow}
\usepackage{hyperref}
\usepackage{framed}
\usepackage{enumitem}
\hypersetup{colorlinks=true,linkcolor=blue!70!black}
\linespread{1.5}

\newcommand{\concept}[2]{%
  \vspace{0.3cm}
  \noindent
  \fcolorbox{blue!50!black}{blue!5}{%
    \begin{minipage}{0.92\textwidth}
    \textbf{\textcolor{blue!60!black}{核心概念} · #1} \\[0.15cm]
    {\small #2}
    \end{minipage}%
  }
  \vspace{0.3cm}
}

\newcommand{\sourcebox}[2]{%
  \vspace{0.3cm}
  \noindent
  \fcolorbox{orange!60!black}{orange!5}{%
    \begin{minipage}{0.92\textwidth}
    \textbf{\textcolor{orange!60!black}{史料阅读}} \\[0.1cm]
    {\small #1} \\[0.15cm]
    {\footnotesize\textcolor{gray}{#2}}
    \end{minipage}%
  }
  \vspace{0.3cm}
}

\newcommand{\method}[1]{%
  \vspace{0.3cm}
  \noindent
  \fcolorbox{green!50!black}{green!3}{%
    \begin{minipage}{0.92\textwidth}
    \textbf{\textcolor{green!50!black}{史学方法}} \\[0.1cm]
    {\small #1}
    \end{minipage}%
  }
  \vspace{0.3cm}
}

\newcommand{\compare}[2]{%
  \vspace{0.3cm}
  \noindent
  \fcolorbox{purple!50!black}{purple!3}{%
    \begin{minipage}{0.92\textwidth}
    \textbf{\textcolor{purple!60!black}{跨文明比较} · #1} \\[0.1cm]
    {\small #2}
    \end{minipage}%
  }
  \vspace{0.3cm}
}

\title{\textcolor{blue!70!black}{\Huge\textbf{中国古代史·下卷}}}
\author{}
\date{}

\begin{document}
\maketitle
\thispagestyle{empty}

\newpage

\section{第二课\quad 春秋战国：旧秩序的解体与新世界的锻造}

公元前771年，犬戎攻破镐京，周幽王死于骊山脚下。平王在诸侯护送下东迁洛邑，中国历史进入了东周时代。此后五百余年，这片土地经历了一场根本性的重构：宗法分封制在内外夹击中逐步瓦解，中央集权的郡县制在战火与变法中艰难成形；世袭贵族的政治垄断被打破，一个以知识和才能而非血统为通行证的官僚阶层开始崛起。到公元前221年秦王政灭六国，一个全新的政治实体——大一统帝国——登上了历史舞台。

这一变迁的幅度和深度在古代世界是罕见的。理解春秋战国，关键不在于记住哪一年发生了哪场战役，而在于把握贯穿其中的三条主线：权力重心如何从天子转移到霸主再转移到国君？经济基础如何从井田制转型为土地私有制？社会等级如何从血缘世袭向功绩流动转变？

\subsection{一、霸主政治：旧秩序的权宜之计}

平王东迁之后，周天子丧失了对诸侯的实际控制。关中根据地丢失，王畿收缩到洛阳周边数百里，经济实力和军事实力都已无法支撑天下共主的地位。诸侯国之间的兼并愈演愈烈，中小国家相继被吞并，只剩十余个大国和若干缓冲小国。

\concept{霸主政治的制度逻辑}{
  霸主政治是周天子权威衰落而新权力中心尚未形成的过渡形态。它的核心特征是"名"与"实"的分离：周天子保持天下共主的法理名号，霸主掌握征伐会盟的实际权力。这种双层结构使得政治博弈必须同时在道义和实力两个层面展开——任何一方如果过早撕破名分的面纱，就会面临来自其他各方的联合反制。
}

在诸大国中率先找到突破口的是齐国。齐桓公（公元前685—前643年在位）任用管仲推行内政改革，同时在外交上打出了"尊王攘夷"的旗号。表面上看，齐国是在替天子讨伐不臣、抵御戎狄入侵；实质上，齐国借天子之名召集诸侯会盟、组织联合军事行动、调解各国纠纷，逐步确立了自己在诸侯中的领导地位。

"尊王攘夷"的高明之处在于它完美地回避了当时最敏感的政治问题——谁敢取代天子？任何企图直接夺位的诸侯都将面临道德上的谴责和军事上的联合讨伐。但如果你宣称自己在维护天子权威、在保护中原文明不受外族侵犯，你反而可以名正言顺地获取权力。管仲的这套操作不仅让齐国成为了春秋第一霸主，还为后世提供了一套可复制的政治策略——九百多年后，曹操"挟天子以令诸侯"，运用的几乎完全相同的逻辑。

齐桓公之后的霸主接力依次是晋文公（城濮之战败楚）、楚庄王（问鼎中原）、吴王阖闾和越王勾践。每一位霸主都重复着类似的路径：通过一场决定性的战争树立威望，然后召集盟会确认领导地位。但霸主政治本身无法解决一个内在矛盾——霸权的维系依赖霸主自身的实力，而没有任何制度可以保证霸主国的实力永不衰退。齐桓公一死，齐国旋即陷入内乱；晋文公之后的晋国虽然维持霸权更久，但最终被三家大夫瓜分。霸主的更替并不导向稳定，它只是延缓了旧秩序崩溃的速度。

\method{如何理解政治制度的"过渡性质"}{
  历史学家分析制度时，常常使用"transitional institution"（过渡性制度）这一概念。它的意思是：某些制度形态存在的目的不是建立一种新的稳定秩序，而是为旧秩序的退出提供缓冲，同时为替代方案的出现争取时间。霸主政治就是典型的过渡性制度——它既不能重建周天子的权威，也无法阻止权力最终流向更集权的君主制，但它在两个世纪的时间内维持了最低限度的国际秩序。

  识别一种制度是"过渡性"的还是"终极性"的，是历史分析中的重要能力。一个简单的方法是追问：如果支撑这一制度的核心条件消失，它是否还能存活？对霸主政治而言，核心条件是大国之间没有一家能够吞并所有其他对手——一旦秦国具备了碾压性的力量优势，霸主政治的"平衡"逻辑就被彻底抛弃了。
}

\subsection{二、铁犁与牛耕：沉默的革命}

如果说霸主政治的更替构成了春秋史的叙事主线，那么真正决定了中国历史走向的力量并不在会盟台上，而在田间地头。公元前6世纪前后，冶铁技术经草原丝绸之路传入中原。铸铁比青铜成本更低、硬度更高，铁制犁铧可以翻耕比石犁和木犁深得多的土层。与此同时，用牛代替人力拉犁在黄河流域逐步推广。据农业史学者推算，铁犁牛耕的组合使单个劳动力的可耕种面积较此前提高约三至五倍。

\begin{center}
\renewcommand{\arraystretch}{1.3}
\begin{tabular}{|c|c|c|c|}
\hline
\textbf{农具形态} & \textbf{材料} & \textbf{动力来源} & \textbf{可耕地/劳力} \\
\hline
耒耜 & 木/石 & 人力 & 约5-10亩 \\
青铜镢 & 青铜 & 人力 & 约10-20亩 \\
铁犁 & 铸铁 & 人力 & 约20-40亩 \\
铁犁+牛耕 & 铸铁 & 畜力 & 约40-60亩 \\
\hline
\end{tabular}
\end{center}

生产效率的大幅提升带来了一个连锁反应。在井田制下，土地被划分为"公田"（产出归领主）和"私田"（产出归耕作者），农民须先耕种公田再耕种私田。当铁犁牛耕使私田的单位产出远超公田时，农民倾向于将更多劳动投入私田，公田逐渐荒废。领主在公田上的收益持续减少，仅仅依靠传统的劳役地租已经无法维持财政。

公元前594年，鲁国实行了一项在财政史上具有里程碑意义的改革——"初税亩"。改革的内容看似简单：不再区分公田和私田，一律按照实际耕种面积征收实物税。但这项政策隐含的制度后果极为深远。它等于承认了土地的实际占有者——不论他是传统上的领主还是实际耕种的农民——对土地产出享有合法支配权。一旦"按亩征税"成为规则，土地私有的事实就得到了制度的默许。

到了战国时期，土地买卖已经成为普遍现象。《韩非子》记载中提到了"中牟之人弃其田耘、卖宅圃而随文学者"——有人卖掉田产和宅地去追逐学问。如果土地不是私有的，买卖就不可能发生。

\concept{井田制到土地私有——一次基础产权制度的更替}{
  从井田制到土地私有制的转变，本质上是一次产权制度的更替。在周代封建体制下，土地的所有权属于周天子（"普天之下，莫非王土"），使用权层层下分，但不能买卖。春秋战国的变革使土地从"不可交易的身份性资产"转变为"可以买卖的财产"。这一转变的意义不亚于英国16世纪的圈地运动——它释放了生产效率，但也制造了新一轮的社会分化和流民问题。
}

\subsection{三、变法的逻辑：从贵族分权到君主集权}

经济基础的变动必然带动上层建筑的调整。到战国时期，各国君主面临的核心问题已经从"如何在国际上树立威信"转变为"如何在激烈的兼并战争中生存下来"。答案是变法——通过制度化的手段强化国家对人口、资源和武力的直接控制。

战国变法的总设计师们虽然散处各国，改革的重心各有侧重，但底层逻辑高度一致：打破旧贵族对土地、兵源和行政资源的垄断，建立一套直通国君的统治体系。魏国的李悝变法最早起步，制定了中国历史上第一部系统的成文法典《法经》；楚国的吴起变法针对贵族世袭特权；韩国的申不害变法侧重君主驾驭臣下的技术。其中最为彻底、影响最为深远的是秦国的商鞅变法。

\concept{商鞅变法为什么在同一时期最有效？}{
  商鞅变法几乎触及了社会的每一个层面。经济上，废井田、开阡陌，承认土地私有并按亩征税，结束了领主对地租的垄断。军事上，建立起等级分明的二十等军功爵位制：任何人都可以凭借战场上斩获的敌人首级晋升爵位，获得相应的土地、房屋和免税特权。一个出身底层的士兵，如果在战场上连续立功，理论上可以一路晋身到贵族阶层——这在分封制下是无法想象的。行政上，在全国推行县制，由国君直接任免行政长官，废除世袭封邑。社会控制上，实行什伍连坐，十家编为一什、五家编为一伍，一家犯法其余各家若不举报则一同受罚。
}

商鞅变法的历史结果所有人都知道：秦国从一个半边缘化的西方诸侯国崛起为战国最强大的军事机器，在不到一个半世纪后完成了统一。但变法的过程绝不是一帆风顺的。商鞅的每一项措施都在得罪旧贵族——军功爵制让失去世袭特权的旧贵族怒火中烧，县制剥夺了他们的封地控制权，什伍连坐则把整个社会变成了互相监视的网络。秦孝公在世时商鞅有最高支持者的庇护，孝公一死，旧贵族的反扑立刻落在了他身上。

公元前338年，秦孝公去世，太子驷继位。被商鞅处罚过的公子虔等人告发商鞅谋反，秦惠文王下令逮捕商鞅。商鞅在逃亡途中试图投宿旅店，旅店主人因他没有官府凭证而拒绝让他入住——这就是商鞅自己制定的法律。最终商鞅被车裂处死，家族被灭。但司马迁在《史记》中写道："秦法未败也。"法律已经嵌入了秦国的制度结构中，制定者死了，制度还在运转。

\concept{制度化和制度化依赖——理解为什么"秦法未败"}{
  商鞅死后秦法并未被废止，这种现象在制度变迁理论中被称为"制度化依赖"。当一项改革已经深入到行政日常（各级官员按法律运作）、社会预期（底层士兵确信斩首可升爵）和经济利益分配（私田主按亩交税而非服劳役）中时，即使最初的推动者消失，既得利益格局也会阻止全面倒退。秦惠文王可以杀商鞅个人，但不可能在不动摇自己统治的前提下废黜整个军功爵制——因为秦国的军事优势就建立在这一制度上。
}

\subsection{四、百家争鸣：观念战场上的秩序重建}

春秋战国不仅是权力、土地和军队的重组期，也是观念世界的秩序重建期。当旧的宗法伦理失去了约束力，各社会阶层都以自己的方式提出了关于"应当如何组织社会"的主张。这就是百家争鸣的实质——它不仅是少数天才思想家的个人创作，更是社会结构转型在观念层面的必然产物。

儒家是这场思想竞赛中影响最为深远的一个流派。孔子（公元前551—前479年）的核心关切是"礼崩乐坏"——他认为社会的混乱源于人们忘记了社会规范。他的解决方案是"仁"（内在的对他人的关爱）和"礼"（外在的行为规范）的统一。孔子的思想在生前并未得到广泛采纳——他周游列国十四年，始终没有遇到一个愿意全面施行其主张的君主。这与春秋末期紧迫的军事和政治需求有关：孔子的方案是"拔本塞源"式的道德改进，见效太慢，无法满足各国君主对迅速富国强兵的渴求。

到了战国时期，孟子（公元前372—前289年）进一步发展了儒家学说。他的中心论点是"性善论"——人天生具有恻隐之心、羞恶之心、辞让之心和是非之心，统治者只需要通过教育把这些"善端"充分扩充出来，就能实现理想的治理。孟子还提出了当时最激进的政治主张："民为贵，社稷次之，君为轻。"在君主权力不断膨胀的战国时代，这个论断显得格外突出。

儒家最大的弱点是体系性的：它始终没有给出一个在竞争环境中快速见效的治理方案。这一弱点被法家敏锐地抓住了。

法家的集大成者是韩非子（约公元前280—前233年）。他师从儒家学者荀子，却得出了与儒家截然相反的结论。韩非认为人性不可改变——人永远追逐利益、规避惩罚。因此治理国家的唯一可靠方法是依靠制度："法"（成文法律）、"术"（君主操控臣下的权术）和"势"（君主的绝对权威）三者的结合。韩非完全否定了儒家"教化治国"的逻辑——如果法律足够严格、赏罚足够分明，就算是一个平庸的君主也能把国家治理好；反之，如果法律松弛，即使最圣明的君主也无能为力。从这个角度看，法家是中国最早的"制度主义"政治哲学。

道家的老子和庄子走了一条截然不同的路。老子在《道德经》中提出"治大国若烹小鲜"——频繁翻动会碎掉，频繁干预会乱掉。最好的治理是"无为"——统治者尽量减少干预，让社会按其自身规律运转。庄子更加彻底：他认为政治本身就是一个不值得投入精力的领域，人生的最高境界是在精神上与万物合一。

墨家（墨子）的立场介于儒家和法家之间。他主张"兼爱"——不分亲疏厚薄地爱一切人——这比儒家的"爱有差等"（先爱家人再推及他人）更为激进。他还主张"非攻"，反对一切侵略战争，并亲自带着弟子们奔走于各国之间阻止战争爆发。墨家最独特的地方在于其组织方式：它是一个纪律严明、成员生活简朴的团体，其最高领袖（钜子）拥有几乎绝对的权威。

\compare{轴心时代——平行的时间窗口，独立的爆发}{
  德国哲学家雅斯贝尔斯在1949年出版的《历史的起源与目标》中注意到一个令人困惑的现象：公元前800年至公元前200年之间，在完全没有交通和传播联系的几个文明区域内，几乎同时涌现出了决定此后数千年精神走向的思想家。中国的孔子和老子、印度的佛陀、伊朗的琐罗亚斯德、以色列的以赛亚诸先知、希腊的苏格拉底和柏拉图——这些名字各自代表着一套根本性的精神突破。雅斯贝尔斯将这一时段称为"轴心时代"。

  轴心时代假说至今仍是思想史上的重要命题，但也受到不少批评。主要批评有三：第一，这些文明之间的"同时性"可能是统计上的偶然——如果将时间窗口放宽，几乎所有文明都能找到某个"思想密集产出"的阶段；第二，轴心期描述的现象主要集中在欧亚大陆，撒哈拉以南非洲和美洲的同期思想发展在雅斯贝尔斯的框架中未被充分讨论；第三，轴心时代的概念可能隐含着一个未经证明的前提——突破发生在特定的"窗口期"且不可重复。

  尽管有这些争议，"轴心时代"作为思考工具是有用的。它迫使我们去追问：为什么许多文明都在社会失序和秩序重建的过程中走向了对根本价值的追问？有没有某些共通的物质条件（文字普及、城市化、跨区域贸易）催生了这种精神劳作？
}

到战国末期，百家争鸣的局面逐渐结束。秦朝统一后"焚书坑儒"，法家成为官方意识形态；汉武帝"独尊儒术"后，儒家重新占据主流。但这并不意味着其他学派的消失——法家的制度工具被后世王朝暗中继承，道家的思想成了失意文人的精神退路，墨家在民间秘密结社中长期延续。百家争鸣的遗产不是一个学派的胜利，而是一整套思想库——后世的中国政治无论怎样演变，都没有超出在这场争论中划定的知识地盘。

\subsection{五、走向统一：秦灭六国的结构性条件}

公元前230年至前221年，秦国在十年之内依次灭亡韩、赵、魏、楚、燕、齐六国。这绝不是一次军事上的偶然爆发——它是一百多年制度竞赛的最终结果。

秦国的优势是系统性的。经济上，铁犁牛耕在秦地的推广配合水利工程（都江堰、郑国渠）使关中平原的粮食产量大幅提升，为长期战争提供了稳定的后勤支撑。军事上，军功爵制确保了士兵的作战动力。行政上，郡县制和高效的文书行政系统使国家动员能力远超六国——云梦睡虎地出土的秦简中包括了详尽的《田律》《仓律》《徭律》，精确规定了每一亩地应交多少粮食、每一仓粮应如何登记、每一年应征发多少劳力服徭役。相比之下，六国普遍保留了更多的封建残余，动员效率远逊于秦国。

\concept{官僚制与非人格化治理——秦国超出六国的制度优势}{
  一个重要的观察：秦国的治理已经开始向"非人格化"的方向移动。官员由中央任免，按照成文法规行使职权，行政档案被严格保存以备稽考。云梦秦简中有一条律文规定，掌管仓库的官员如丢失物品，必须在规定期限内赔偿，否则将面临严厉处罚。这种"按制度追责"的模式在六国中并不多见。制度史学者认为，秦国统一的一个重要条件不在于它"更强"，而在于它建立了一套可重复的、可监控的、不依赖于个别官员品德的高效行政系统。
}

\method{单一变量分析 vs 系统分析——哪种解释更可靠？}{
  许多通俗史书喜欢用某一个原因来解释秦的统一：有人说是商鞅变法，有人说是长平之战消灭了赵国主力，有人说是地理优势（关中易守难攻）。这些解释都有一定道理，但单独拿出来都无法令人满意。

  历史解释中的一个基本方法论要求是：你提出的原因必须足以解释结果的大小。如果商鞅变法是秦统一的主要原因，那么为什么变法后一百多年秦国才完成统一？如果长平之战是关键，那么为什么赵国覆灭后秦国又用了十九年才灭掉其余五国？单一变量解释的缺陷在于它把复杂系统简化为线性的因果链条。

  更可靠的分析方式是系统分析：把秦国的农业产出、军事动员、行政效率、外交策略和对手的对应条件放在同一个框架内进行综合评估。在这些因素中，哪一个单独拿出来可能都不足以"决定"统一，但它们叠加在一起形成的系统优势，在长期竞争中几乎是不可逆的。
}

公元前221年，秦王政在咸阳登基称帝，自称"始皇帝"。他废除了延续近千年的分封制，将全国划分为三十六郡，由中央直接派遣郡守和县令。他统一了文字（小篆和隶书）、车轨宽度、度量衡和货币。他把六国的兵器收缴到咸阳熔化，铸成十二尊铜人立在宫殿前。他拆毁了六国边境上的关塞堡垒，修筑了通往全国的驰道网络。

这些措施的目的只有一个：确保这片土地再也不会回到战国群雄割据的局面。分裂持续了五百多年，而秦朝的寿命只有十五年。但秦朝塑造的制度框架——郡县制、官僚制、"书同文、车同轨"——被后世继承了两千年。这个矛盾本身就给出了关于春秋战国这段历史最深刻的评价：在血与火中孕育出来的制度胚胎，远远比制造它的那个政权活得更久。

\subsection*{A组：内容复述与要点整理}

\begin{enumerate}[label=A\arabic*.]
  \item 简述单主政治的形成原因和运作特征。为什么齐桓公需要通过"尊王攘夷"而不是直接取代周天子来获取权力？
  \item 铁犁牛耕的普及如何导致了井田制的解体？试按照"技术改进→劳动效率变化→土地制度调整"的逻辑链条做出解释。
  \item 商鞅变法的主要内容有四个方面。请分别说明每一项改革针对的是什么问题，并解释二十等军功爵制为什么能够改变秦国的社会结构。
  \item 百家争鸣中哪四个学派最为重要？列出每个学派的核心理念，并说明它们在"如何治理社会"这一根本问题上的分歧。
  \item 秦国的统一优势涉及经济、军事、行政和外交四个层面。请各自举出一个具体例证。
\end{enumerate}

\subsection*{B组：深入分析与学术写作}

\begin{enumerate}[label=B\arabic*.]
  \item \textbf{概念应用题}：本课正文引入了一个分析工具——"过渡性制度"。请用这一概念分析以下任意一个历史案例（约400字）：(a) 汉武帝的推恩令，(b) 唐代的藩镇制度，(c) 清末的预备立宪。你需要说明：该制度是在什么旧的制度框架崩溃时出现的？它为什么不能作为长期稳定的解决方案？

  \item \textbf{史料分析题}：阅读以下两段史料后回答：
  \begin{quote}
    \textbf{史料一——司马迁《史记·商君列传》：} 商君相秦十年，宗室贵戚多怨望者。后五月而秦孝公卒，太子立。公子虔之徒告商君欲反，发吏捕商君。秦惠王车裂商君以徇，曰："莫如商鞅反者！"遂灭商君之家。
  \end{quote}
  \begin{quote}
    \textbf{史料二——《战国策·秦策一》：} 商君治秦，法令至行，公平无私，罚不讳强大，赏不私亲近，法及太子，黥劓其傅。期年之后，道不拾遗，民不妄取，兵革大强，诸侯畏惧。
  \end{quote}
  \begin{enumerate}[label=(\alph*)]
    \item 史料一和史料二对商鞅变法的评价有何差异？分别倾向于从什么视角（改革者个人命运 vs 改革的社会效果）做出评判？
    \item 司马迁说"秦法未败也"——结合这两段史料，你认为"法未败"的判断建立在什么证据之上？商鞅的个人遭遇和秦法的延续之间是否存在矛盾？
  \end{enumerate}

  \item \textbf{论证题}（600—800字）：韩非子认为"不期修古，不法常可"，治理国家不应效法古代、不应遵循一成不变的成规。请结合春秋战国时期至少两个具体案例，论证或反驳这一命题。要求：(a) 明确提出你的核心论点，(b) 每个案例的描述控制在100字以内，(c) 在论证中加入对"如果……可能……"的反事实推理。

  \item \textbf{比较研究题}：回顾上卷第一课中关于周代分封制与欧洲封建制度的对比，结合本课关于春秋战国社会转型的论述，建立一份不少于600字的比较分析框架，回答：中国为什么从分裂走向了中央集权的统一帝国，而欧洲在中世纪之后却走向了多国并立的民族国家体系？你需要考虑以下变量：地理条件（平原 vs 破碎地形）、继承制度（嫡长子 vs 诸子平分）、文字统一程度、外部威胁的性质（草原游牧 vs 多方敌人）和思想传统（儒家的大一统意识 vs 多元权力结构）。
  \item \textbf{研究性学习打包：}
  \begin{enumerate}[label=(\alph*)]
    \item 查询你所在城市或省份在春秋战国时期所属的诸侯国。阅读该诸侯国在《史记》中的相关"世家"，撰写主题为"一个中等诸侯国的生存策略"的简章（800-1000字）。你需要讨论该国的地理位置、外交选择和最终的命运。
    \item \textbf{对比阅读（可选）：}对照《春秋》和《左传》对同一事件的记载（例如"郑伯克段于鄢"在两种文本中的表述差异）。分析编年体史书和叙事体史书在信息选择、详略安排和隐含评价上的区别。将你的发现写成一份分析笔记。
  \end{enumerate}
\end{enumerate}

\subsection*{参考书目（分阶难度标注）}

\begin{itemize}
  \item [★] 李开元《秦崩：从秦始皇到刘邦》，生活·读书·新知三联书店
  \item [★] 杨宽《战国史》（青少年版），上海人民出版社
  \item [★★] 钱穆《国史大纲》第一编第三至五章，商务印书馆
  \item [★★] 黄仁宇《中国大历史》第三章，生活·读书·新知三联书店
  \item [★★★] 许倬云《中国古代社会史论》，广西师范大学出版社
  \item [★★★] 雅斯贝尔斯《历史的起源与目标》，华夏出版社
\end{itemize}

\end{document}
